\begin{ouverture}[geometrie]{17}{La symétrie axiale}{Partie 4 · Espace et géométrie}
{Axes · Pliage · Quadrillage · Figures à compléter}
\objectif{Reconnaître les axes de symétrie d’une figure.}
\objectif{Compléter une figure pour la rendre symétrique.}
\objectif{Construire un symétrique sur quadrillage.}
\prerequis{Quadrillages, perpendiculaires et distances.}
\end{ouverture}
\pagestyle{otMain}\essentiel\fontsize{11}{15}\selectfont
\phantomsection\label{cours-17-1}
\essSection{1}{Reconnaître une symétrie}
\begin{definition}Une droite est un axe de symétrie d’une figure si, en pliant la feuille le long de cette droite, les deux parties de la figure se superposent exactement. Le calque, retourné le long de l’axe, fait la même vérification.\end{definition}
\label{fig-cours-17-1-cours-symetrie}\figureGeo{cours-symetrie}

\begin{exemple}Le carré a 4 axes de symétrie, le rectangle en a 2, le triangle équilatéral 3, le cercle une infinité. Certaines figures n’en ont aucun : un pliage qui déborde suffit à l’affirmer.\end{exemple}
\begin{attention}Les diagonales du rectangle ne sont pas des axes de symétrie : le pliage le long d’une diagonale ne superpose pas les deux moitiés. Toujours imaginer ou réaliser le pliage.\end{attention}
\clearpage
\phantomsection\label{cours-17-2}
\essSection{2}{Construire sur quadrillage}
\begin{savoirfaire}[Compter les carreaux]Sur quadrillage, avec un axe vertical : le symétrique d’un point se place de l’autre côté, à autant de carreaux de l’axe, sur la même ligne. Avec un axe horizontal, on compte les carreaux au-dessus et au-dessous.\end{savoirfaire}
\begin{propriete}La symétrie conserve les longueurs et les formes : le symétrique d’un segment de 4 carreaux mesure 4 carreaux. Un point situé sur l’axe est son propre symétrique : il ne bouge pas.\end{propriete}
\begin{savoirfaire}[Compléter une figure]Pour compléter une figure dont l’axe est tracé : traiter les sommets un par un, placer chaque symétrique, puis relier dans le même ordre. Contrôler en pliant mentalement.\end{savoirfaire}
\begin{exemple}Les lettres offrent de bons exemples : A possède un axe vertical, B un axe horizontal, H les deux, et F n’en a aucun.\end{exemple}

\clearpage\rubrique{Découverte}
\begin{blocLivre}{ex-17-1}\begin{exercice}[1]{Axe ou pas ?}
Pour chaque figure dessinée, dire si la droite tracée en pointillés est un axe de symétrie.
\label{fig-ex-17-1}\figureGeo{sym-1}
\end{exercice}\end{blocLivre}
\begin{blocLivre}{ex-17-2}\begin{exercice}[1]{Compter les axes}
Donner le nombre d’axes de symétrie : du carré, du rectangle et du triangle équilatéral dessinés.
\label{fig-ex-17-2}\figureGeo{sym-2}
\end{exercice}\end{blocLivre}
\begin{blocLivre}{ex-17-3}\begin{exercice}[1]{La diagonale piégée}
Milo plie un rectangle le long d’une diagonale et dit : « C’est un axe de symétrie, les deux triangles sont pareils. » Que lui répondre ?
\label{fig-ex-17-3}\figureGeo{sym-3}
\end{exercice}\end{blocLivre}
\begin{blocLivre}{ex-17-4}\begin{exercice}[1]{Le symétrique d’un point}
Placer le symétrique de chaque point marqué par rapport à l’axe vertical tracé.
\label{fig-ex-17-4}\figureGeo{sym-4}
\end{exercice}\end{blocLivre}
\begin{blocLivre}{ex-17-5}\begin{exercice}[1]{Compléter la maison}
Compléter le dessin de la maison pour que l’axe vertical soit un axe de symétrie.
\label{fig-ex-17-5}\figureGeo{sym-5}
\end{exercice}\end{blocLivre}

\rubrique{Entraînement}
\begin{blocLivre}{ex-17-6}\begin{exercice}[2]{Axe horizontal}
Construire le symétrique du segment [AB] par rapport à l’axe horizontal tracé.
\label{fig-ex-17-6}\figureGeo{sym-6}
\end{exercice}\end{blocLivre}
\begin{blocLivre}{ex-17-7}\begin{exercice}[2]{Le papillon}
Compléter le papillon pour que la droite horizontale soit un axe de symétrie.
\label{fig-ex-17-7}\figureGeo{sym-7}
\end{exercice}\end{blocLivre}
\begin{blocLivre}{ex-17-8}\begin{exercice}[2]{Les lettres}
Parmi les lettres A, B, H et F dessinées, dire lesquelles ont un axe de symétrie vertical, horizontal, ou aucun.
\label{fig-ex-17-8}\figureGeo{sym-8}
\end{exercice}\end{blocLivre}
\begin{blocLivre}{ex-17-9}\begin{exercice}[2]{Chercher l’erreur}
Léa a construit le symétrique du triangle par rapport à l’axe vertical, mais un sommet est mal placé. Lequel ? Corriger.
\label{fig-ex-17-9}\figureGeo{sym-9}
\end{exercice}\end{blocLivre}
\begin{blocLivre}{ex-17-10}\begin{exercice}[2]{Sur l’axe}
Le segment [AB] a une extrémité A posée sur l’axe. Construire son symétrique. Que devient le point A ?
\label{fig-ex-17-10}\figureGeo{sym-10}
\end{exercice}\end{blocLivre}

\rubrique{Problèmes}
\begin{blocLivre}{ex-17-11}\begin{exercice}[3]{Le carrelage}
Compléter le motif du carrelage pour que l’axe vertical soit un axe de symétrie, en coloriant les bonnes cases.
\label{fig-ex-17-11}\figureGeo{sym-11}
\end{exercice}\end{blocLivre}
\begin{blocLivre}{ex-17-12}\begin{exercice}[3]{Deux axes à la fois}
Compléter la figure pour qu’elle admette l’axe vertical et l’axe horizontal tracés comme axes de symétrie.
\label{fig-ex-17-12}\figureGeo{sym-12}
\end{exercice}\end{blocLivre}
\begin{blocLivre}{ex-17-13}\begin{exercice}[3]{Les longueurs conservées}
Le segment [CD] mesure 4 carreaux. Construire son symétrique par rapport à l’axe, puis donner sa longueur sans la mesurer. Justifier.
\label{fig-ex-17-13}\figureGeo{sym-13}
\end{exercice}\end{blocLivre}
\begin{blocLivre}{ex-17-14}\begin{exercice}[3]{Le bateau du port}
Compléter le dessin du bateau et de sa voile pour que l’axe vertical soit un axe de symétrie du dessin complet.
\label{fig-ex-17-14}\figureGeo{sym-14}
\end{exercice}\end{blocLivre}
\begin{blocLivre}{ex-17-15}\begin{exercice}[3]{Vrai ou faux ?}
Pour chaque phrase, dire si elle est vraie et corriger si besoin : « Le symétrique d’un point est toujours de l’autre côté de l’axe. » ; « Une figure a au plus un axe de symétrie. »
\label{fig-ex-17-15}\figureGeo{sym-15}
\end{exercice}\end{blocLivre}

\clearpage
